\documentclass[12pt]{amsart}
%prepared in AMSLaTeX, under LaTeX2e
\addtolength{\oddsidemargin}{-.65in} 
\addtolength{\evensidemargin}{-.65in}
\addtolength{\topmargin}{-.55in}
\addtolength{\textwidth}{1.3in}
\addtolength{\textheight}{.9in}

\renewcommand{\baselinestretch}{1.05}

\usepackage{verbatim} % for "comment" environment

\usepackage{palatino}

\newtheorem*{thm}{Theorem}
\newtheorem*{defn}{Definition}
\newtheorem*{example}{Example}
\newtheorem*{problem}{Problem}
\newtheorem*{remark}{Remark}

\usepackage{fancyvrb,xspace,dsfont}

\usepackage[final]{graphicx}

% macros
\usepackage{amssymb}
\usepackage[dvipsnames]{xcolor}

\usepackage{hyperref}
\hypersetup{pdfauthor={Ed Bueler},
            pdfcreator={pdflatex},
            colorlinks=true,
            citecolor=blue,
            linkcolor=red,
            urlcolor=blue,
            }

\newcommand{\br}{\mathbf{r}}
\newcommand{\bv}{\mathbf{v}}
\newcommand{\bx}{\mathbf{x}}
\newcommand{\by}{\mathbf{y}}

\newcommand{\cB}{\mathcal{B}}
\newcommand{\cD}{\mathcal{D}}
\newcommand{\cF}{\mathcal{F}}
\newcommand{\cH}{\mathcal{H}}
\newcommand{\cL}{\mathcal{L}}
\newcommand{\cV}{\mathcal{V}}
\newcommand{\cW}{\mathcal{W}}

\newcommand{\CC}{\mathbb{C}}
\newcommand{\NN}{\mathbb{N}}
\newcommand{\RR}{\mathbb{R}}
\newcommand{\ZZ}{\mathbb{Z}}

\renewcommand{\Im}{\mathrm{Im}}
\renewcommand{\Re}{\mathrm{Re}}

\newcommand{\eps}{\epsilon}
\newcommand{\grad}{\nabla}
\newcommand{\lam}{\lambda}
\newcommand{\lap}{\triangle}

\newcommand{\ip}[2]{\ensuremath{\left<#1,#2\right>}}

\newcommand{\image}{\operatorname{im}}
\newcommand{\onull}{\operatorname{null}}
\newcommand{\rank}{\operatorname{rank}}
\newcommand{\range}{\operatorname{range}}
\newcommand{\trace}{\operatorname{tr}}
\newcommand{\Span}{\operatorname{span}}

\newcommand{\prob}[1]{\bigskip\noindent\textbf{#1.}\quad }

\newcommand{\pts}[1]{(\emph{#1 pts}) }
\newcommand{\epart}[1]{\medskip\noindent\textbf{(#1)}\quad }
\newcommand{\ppart}[1]{\,\textbf{(#1)}\quad }

\newcommand{\ds}{\displaystyle}

\newcommand{\nex}{\medskip\noindent}


\begin{document}
\scriptsize \noindent Math 617 Functional Analysis (Bueler) \hfill \emph{assigned January 2026}
\normalsize\medskip

\Large\centerline{\textbf{Assignment 2}}
\large
\medskip

\centerline{\textbf{Due Monday 2 February at the beginning of class}}
\medskip
\normalsize

\thispagestyle{empty}

\bigskip
\noindent This Assignment is based on Chapters 1 and 2 of our textbook,\footnote{K.~Saxe,~\emph{Beginning Functional Analysis}, Springer 2010.} and on the lectures.


\newcommand{\augment}[1]{\quad\,\strut \begin{minipage}[t]{0.7\textwidth} \emph{#1}\end{minipage}}

\bigskip
\noindent \textsc{Do the following} Exercises from Chapter 2 (see pages 29--31):

\begin{itemize}
\item \textbf{Exercise 2.1.1} \augment{Note that Saxe defines \emph{open} and \emph{closed} sets on page 16.  Please use those definitions.}
\item \textbf{Exercise 2.1.4}
\item \textbf{Exercise 2.1.6}
\item \textbf{Exercise 2.1.7} \augment{State your claim, and then prove it.}
\item \textbf{Exercise 2.1.10}
\item \textbf{Exercise 2.1.13} \augment{Do parts \textbf{(a)} and \textbf{(b)} only.}
\item \textbf{Exercise 2.2.2} \augment{You may use, without proof, the fact that a countable union of countable or finite sets is countable.}
\item \textbf{Exercise 2.3.1}
\item \textbf{Exercise 2.3.5} \augment{To show that the formula defines a norm, you may use, without proof, the triangle inequality for integrals, that is, $|\int_E f(t)\,dt| \le \int_E |f(t)|\,dt$.}
\item \textbf{Exercise 2.3.8} \augment{You may use, without proof, the fact that $\RR$ is complete.}
\end{itemize}


\medskip
\noindent \textsc{Do the following additional problems.}

\prob{P3}  \emph{In lecture\footnote{See these slides: \url{https://bueler.github.io/fa/assets/slides/S26/week2.pdf}.} I solved a classic boundary value problem (BVP) by writing $u(x)$ as an integral of the data $\phi(x)$.  This is nearly the same, or very slightly easier.  In fact $\phi \in L^1(0,1)$ suffices, but we need not worry about generality for now.}

\medskip \noindent Consider the following BVP, with data $\phi\in C([0,1])$:
\begin{align*}
-u''(x) &= \phi(x) \quad \text{ for all } x\in (0,1) \\
u(0) &= 0 \\
u'(1) &= 0
\end{align*}
Solve this problem by integration.  In particular, write the solution as
	$$u(x) = \int_0^1 k(x,t) \phi(t)\,dt,$$
and give the precise formula for $k(x,t)$.  Is $k$ a continuous function of its two variables?  Is it bounded?

\end{document}
